%MIT OpenCourseWare: https://ocw.mit.edu
%RES.18-012 Algebra II Student Notes, Spring 2022
%License: Creative Commons BY-NC-SA 
%For information about citing these materials or our Terms of Use, visit: https://ocw.mit.edu/terms.

% \rhead{\rightmark}
\lhead{Lecture 23: Noetherian Rings}

\section{Noetherian Rings}

\vspace{0.75em}

\begin{definition}
    A ring $R$ is \textbf{Noetherian} if every ideal in $R$ is finitely generated. 
\end{definition}

In other words, every quotient can be obtained by imposing a finite number of relations. 
\subsection{Submodules over Noetherian Rings}

The first important property of Noetherian rings is the following:

\begin{proposition}\label{noetheriansubmodules}
    A ring $R$ is Noetherian if and only if every submodule in a finitely generated $R$-module is itself finitely generated. 
\end{proposition}

\begin{corollary}\label{finpresented}
    If $R$ is Noetherian, then every finitely generated module is finitely presented. 
\end{corollary}

Before we prove this, we'll look at a few observations. 

\begin{lemma} \label{noetherianhomomorphisms}
    If we have a surjective homomorphism $\varphi : M \to N$ of $R$-modules, then:
    \begin{enumerate}
        \item If $M$ is finitely generated, then $N$ is also finitely generated. 
        \item If $N$ is finitely generated, and $K = \ker(\varphi)$ is also finitely generated, then $M$ is also finitely generated. 
    \end{enumerate}
\end{lemma}

\begin{note}
    This is one place where the intuition from linear algebra is useful: it's helpful to think about the case where $R$ is a field, so being finitely generated is equivalent to being finite-dimensional. In this case, we know more precisely that \[\dim(M) = \dim(N) + \dim(K).\] We won't get information this precise in the general case, but the proof of finiteness is similar. 
\end{note}

\begin{proof}[Proof of Lemma \ref{noetherianhomomorphisms}]
    The first part is obvious --- take any set of generators $m_1$, \ldots, $m_n$ for $M$. Then their images $\varphi(m_1)$, \ldots, $\varphi(m_n)$ must generate $N$. 

    For the second part, let $k_1$, \ldots, $k_a$ be a set of generators for $K$, and $n_1$, \ldots, $n_b$ a set of generators for $N$. Now pick $\tilde{n}_1$, \ldots, $\tilde{n}_b$ such that $\varphi(\tilde{n}_i) = n_i$ (which we can do by surjectivity). 

    Now we claim that the $k_i$ and $\tilde{n}_i$ generate $M$: given any $x \in M$, we can find $r_1$, \ldots, $r_b$ in $R$ such that \[\varphi(x) = r_1n_1 + \cdots + r_bn_b\] (since the $n_i$ generate $N$). So then we have \[\varphi(x - r_1\tilde{n}_1 - \cdots - r_b\tilde{n}_b) = \varphi(x) - r_1n_1 - \cdots - r_bn_b = 0,\] which means $x - \sum r_i\tilde{n}_i$ is in $K$. So we can express $x - \sum r_i\tilde{n}_i = \sum s_jk_j$, which gives an expression for $x$ as a linear combination of the $\tilde{n}_i$ and $k_j$. 
\end{proof}

\begin{proof}[Proof of Proposition \ref{noetheriansubmodules}]
    One direction is clear: an ideal of $R$, by definition, is the same as a submodule of the free module (which is generated by one element). So if every submodule of a finitely generated module is finitely generated, then $R$ must be Noetherian. 

    For the other direction, we want to show that if every ideal is finitely generated, then so is every submodule of a finitely generated module. 

    The strategy is to reduce to the case of a free module $R^n$, and use induction on $n$ --- we know this is true for $n = 1$, so we want to reduce to this case. 

    Let $M$ be a finitely generated module. Then by picking a set of generators, we can find a surjective homomorphism $\varphi : R^n \to M$ (fixing such a homomorphism is equivalent to fixing a set of generators). 

    Then by the correspondence theorem and the above lemma, it's enough to check that every submodule of $R^n$ is finitely generated (since the submodules of $M$ are exactly the images of the submodules of $R^n$ containing $\ker \varphi$). 

    Now we can argue by induction on $n$. First, the base case $n = 1$ follows directly from the definition of  a Noetherian ring, since submodules of $R$ are exactly ideals. 

    For the inductive step, consider a submodule $N \subset R^n$, with $n > 1$. Now split $R^n = R \times R^{n - 1}$, and take the projection homomorphism $\pi : R^n \to R^{n - 1}$, which sends $(r_1, \ldots, r_n) \mapsto (r_2, \ldots, r_n)$. 

    Then $\pi(N)$ is a submodule of $R^{n - 1}$, so by the induction assumption, it's finitely generated. Meanwhile, the kernel $K$ of $\pi$ is the set of elements of the form $(r, 0, 0, \ldots)$ which are in $N$. But this is a submodule of the free rank-$1$ module $R$, so $K$ is also finitely generated. 

    So by the lemma, since $K$ and $\pi(N)$ are both finitely generated, $N$ is finitely generated as well. 
\end{proof}

\begin{example}
    When considering submodules $N \subset \mathbb{Z}^2$, we'd take the points on the $x$-axis as $K$, and the projections onto the $y$-axis as $\pi(N)$. Note that $N$ is not necessarily $K \times \pi(N)$. 
\end{example} 

\begin{question}
Did we use the fact that finitely generated modules are finitely presented in this proof, when we obtained the homomorphism $R^n \to M$?
\end{question}

\begin{ans}
No --- the surjective homomorphism $\varphi : R^n \to M$ comes just from the definition of being finitely generated. We take some generators $m_1$, \ldots, $m_n$, and then map $(r_1, \ldots, r_n)$ to the linear combination $r_1m_1 + \cdots + r_nm_n$. Being finitely presented would only matter if we were looking at the \emph{kernel} of this homomorphism (which we didn't need to do here). 
\end{ans}

So now we have that if $R$ is Noetherian, any finitely generated module is also finitely presented:

\begin{proof}[Proof of Corollary \ref{finpresented}]
    If the module is finitely generated, then there is a surjective map $\varphi : R^n \to M$. Then $\ker\varphi$ is a submodule of $R^n$, so it must be finitely generated as well. 
\end{proof}

This means the classification of finitely presented abelian groups that we saw earlier is actually a classification of all finitely \emph{generated} abelian groups. 

\begin{note}
It's also possible to define Noetherian rings by Proposition \ref{noetheriansubmodules} (as rings with the property that every submodule of a finitely generated module is finitely generated). But we gave the definition using ideals so that the property that's easier to check is in the definition, and the one that's more useful is in the proposition. 
\end{note}

\subsection{Constructing Noetherian Rings}

So we've seen why the notion of Noetherian rings is useful. But in order to use it, we want to see how to produce more examples of Noetherian rings, beyond just fields and PIDs. 

First, there is a simple observation we can make:

\begin{lemma}
    A quotient of a Noetherian ring is again Noetherian --- if $R$ is a Noetherian ring and $I$ an ideal of $R$, then the ring $S = R/I$ is also Noetherian. 
\end{lemma}

\begin{proof}
    This is immediate from the correspondence theorem: an ideal in $R/I$ is of the form $J/I$, where $J \subset R$ is an ideal containing $I$. Then we can just take the images of the generators --- if $J = (x_1, \ldots, x_n)$, then $J/I = (\overline{x}_1, \ldots, \overline{x}_n)$. So all ideals of $R/I$ are finitely generated. 
\end{proof}

\begin{note}
    A subring in a Noetherian ring is not necessarily Noetherian --- so this concept is more subtle than the dimension of a vector space. 

    For example, $\mathbb{C}[x, y]$ is Noetherian. But the subring $\mathbb{C} + x\mathbb{C}[x, y]$ (consisting of polynomials which are constant mod $x$) is \emph{not} Noetherian --- the ideal $x\mathbb{C}[x, y]$ is not finitely generated. 
\end{note}

\subsubsection{Hilbert Basis Theorem}

There's actually a powerful tool that shows many rings are Noetherian: 

\begin{theorem}[Hilbert Basis Theorem] \label{hilbertbasis}
    If $R$ is Noetherian, then $R[x]$ is also Noetherian. 
\end{theorem}

% \begin{remark}
%     When Hilbert proved this theorem in 1890, there's a legend that a famous mathematician Paul Gordan (referred to as the king of invariant theory) said this is not mathematics, it's theology. Previously, people had to work with rings case by case, and concretely produce finitely many elements generating an ideal. In contrast, this theorem has a very abstract proof that doesn't give much information about how to actually write down the generators. 
% \end{remark}

This theorem has useful implications:

\begin{corollary}
    If $R$ is Noetherian, then $R[x_1, \ldots, x_n]/I$ is also Noetherian, for any ideal $I$. 
\end{corollary}

So if we start with a field, this shows us how to produce many examples of Noetherian rings. 

\begin{corollary}
    Any algebraic subset in $\mathbb{C}^n$ --- a subset given by a collection of polynomial equations --- is always given by a \emph{finite} set of polynomial equations. 
\end{corollary} 

\begin{proof}[Proof of Theorem \ref{hilbertbasis}]
    Let $I \subset R[x]$ be an ideal, so we want to check that $I$ is finitely generated. It's enough to find a finite collection of polynomials $P_1$, \ldots, $P_n$ in $I$ and a bound $d$, such that every element in $I$ can be reduced to a polynomial of degree $d$ using the $P_i$ --- meaning that \[I \subset (P_1, \ldots, P_n) + R[x]_{\leq d}.\] In other words, once the polynomials and $d$ are fixed, then for every $P \in I$, we need to be able to find $Q_1$, \ldots, $Q_n$ in $R[x]$ such that \[\deg(P - \sum Q_iP_i) \leq d.\] 
    
    If we can find such $P_i$ and $d$, then \[I \subset (P_1, \ldots, P_n) + (I \cap R[x]_{\leq d}).\] But the second term is finitely generated over $R$ (it's not a submodule of $R[x]$, but it \emph{is} a submodule of $R$), since $R[x]_{\leq d}$ is a free module of rank $d + 1$ over $R$, and $R$ is Noetherian. So if it's generated by $S_1$, \ldots, $S_m$, then $I$ is generated by the $P_i$ and $S_i$. 

    So now we want to figure out how to do this --- in some sense, the idea is generalized division with remainder. 

    Consider the ideal $\overline{I}$ in $R$ consisting of the \emph{leading coefficients} of polynomials in $I$ (along with $0$) --- we saw in the homework that this is an ideal. Then $\overline{I}$ is finitely generated. Let $P_1$, \ldots, $P_n$ be polynomials whose leading coefficients generate $\overline{I}$, and let $d = \max(\deg P_i)$. 

    Now if we have a polynomial $P$ of degree greater than $d$, we can cancel its leading coefficient --- we can find $Q_i$ such that $\sum Q_iP_i$ has the same degree and same leading coefficient as $P$, and then subtract them to decrease the degree of $P$ by at least $1$. We can then repeat this until $P$ has degree at most $d$. 
\end{proof}

% \begin{note}
%     This proof is not very constructive --- it's not an effective way of finding generators. 
% \end{note}

\begin{question}
Why is $I \cap R[x]_{\leq d}$ finitely generated?
\end{question}

\begin{ans}
The point is that $R[x]_{\leq d}$ is a free module of rank $d + 1$ over $R$ --- we forget that we can multiply the polynomials and just add them and scale by elements of $R$, so the coordinates correspond to the coefficients of $1$, $x$, \ldots, $x^d$. Then $I \cap R[x]_{\leq d}$ is a submodule of $R[x]_{\leq d}$, which is a finitely generated module over $R$; so since $R$ is Noetherian, $I \cap R[x]_{\leq d}$ is also finitely generated. 
\end{ans}

\subsection{Chain Conditions}

\vspace{0.75em}

\begin{proposition} \label{chaincondition}
    A ring is Noetherian if and only if every increasing chain of ideals stabilizes. In other words, if there is a chain of ideals $I_1 \subseteq I_2 \subseteq \cdots$, then from some point on, $I_n = I_{n + 1} = \cdots$. 
\end{proposition}

In other words, $R$ is \emph{not} Noetherian if and only if there exists an infinite chain $I_1 \subsetneq I_2 \subsetneq \cdots$ of ideals. 

\begin{proof}[Proof Outline]
If $R$ is Noetherian and we have a chain $I_1 \subseteq I_2 \subseteq \cdots$, then their union $I = I_1 \cup I_2 \cup \cdots$ is an ideal. Since $R$ is Noetherian, then $I$ is finitely generated. So each of the generators must have some ideal $I_k$ that it first shows up in, and since there's finitely many, there's some ideal $I_k$ containing all the generators of $I$ (which must equal $I$). 

For the other direction, we can essentially take an ideal $I$ that isn't finitely generated, and starting with the empty ideal, we keep adding an element of $I$ which isn't in any of the previous ideals. 
\end{proof}
